\section{讨论}
\label{sec:discussion}

\subsection{预算分档的适用边界}

预算分档编排的收益取决于任务分布、基础模型能力和成本约束之间的匹配关系。MBPP 题面更杂、单步残差更大，修补、分流救援和级联能够逐级抬高通过率；HumanEval 已接近饱和，后续档位主要收束少量残差题。两组结果共同说明，预算分档的价值来自按任务动态范围配置计算投向。

级联档的定位也更清晰。它在 MBPP 上提供最高通过率，同时消耗更多令牌；在 HumanEval 或部分模型上，级联相对中档的增益较小。实际部署中，级联适合高价值样本、离线批处理或追求上限的场景；在线常规场景则更适合以单步+修补和分流主链作为默认路径。

\subsection{机制证据的解释}

学习式路由基线显示，当前静态结构特征对运行失败标签的直接判别能力接近随机；这并不削弱结构分的作用，反而说明它更适合作为预算排序信号。结构分把题面复杂、接口复杂或测试条件密集的样本优先送入重路径，端到端前沿、阈值消融和固定多重集对照共同支撑了这一用法。

固定多重集随机路由对照进一步拆开了覆盖率与结构对齐。该实验固定轻/重路径数量，只改变题目与槽位的对应关系。MBPP 上结构路由多通过 4 题但消耗更多令牌，说明结构对齐带来可见收益，也对应更高计算投入；HumanEval 上通过数相同，则与其饱和分布一致。

反馈形态对照给出了明确的工程结论。五随机种子结果中，压缩回注与完整回溯取得相近通过率，且平均令牌更低。压缩回注保留错误类型、关键片段和定向修补提示，减少了无关运行栈文本，是本文中预算档位的默认反馈形式。

\subsection{部署建议}

基于主结果与消融，预算策略可以分为两层。中预算场景下，优先保留单步+修补、结构化主链与压缩失败回注，这些模块在 MBPP 上具有较好的成本--收益关系，也更容易解释：它们把失败信息转化为结构化约束，帮助模型完成局部修补或更聚焦的重链生成。

高预算场景下，新增计算应优先投入候选扩展，再视成本约束开启精修。级联消融显示，候选扩展是高预算档的核心增益来源；精修相关模块更多承担细粒度折中。若预算只允许增加一个高档模块，候选扩展通常比后处理精修更值得保留。

落地前仍需在目标业务数据上做小规模回归。简单模板题会削弱级联收益，复杂边界条件、长测试或多约束签名会放大分流与候选扩展的价值。预算分档协议的优势在于可以用同一记账口径快速完成这种回归。

\subsection{局限与未来工作}

本文仍有若干局限。细粒度消融以 MBPP 为主；HumanEval 更适合作为残差收束观察。多模型表与多数消融表采用随机种子 42，用于展示形态和机制趋势；固定多重集路由对照目前也是单次配对，后续可扩展为多随机种子实验。

此外，本文将总令牌作为接口成本主代理，真实部署还会考虑按次计费、延迟、并发限制与失败重试策略。未来工作将把 $\rho$ 扩展为包含固定调用成本、延迟惩罚和预算上限的联合指标，并持续跟踪模型版本更新带来的单步能力、反馈敏感性和候选多样性变化。
